\scanpage{90}\mspage{34}
% L1-2
pour les foncteurs induits sur les loc. par les
\[
 \underline{\mathrm{Fibsc}}\,S\longrightarrow\underline{\mathrm{Fib}}\,S,\quad
 \underline{\mathrm{Cofibsc}}\,S\longrightarrow\underline{\mathrm{Cofib}}\,S.
\]
% L3-7
Comme les localisées des catégories buts des ces deux foncteurs ont tout l'air de s'identifier aux sous-catégories précédentes \deleted{\unsure{d}}d'objets loc constants, il semblerait donc que \textcircled{\scriptsize 2\textsuperscript{o}} les foncteurs précédents définissent
% L8-11
des rétractions de la catégorie $\mathrm{HOT}_W(S)$ resp. $\mathrm{HOT}_W(S^{\unsure{\mathrm o|\circ}})^{\mathrm o}$ sur la sous-catégorie des objets loc. constants sur $S$, les foncteurs localisés des foncteurs d'inclusion
% L12
\[
 \underline{\mathrm{Fibsc}}_0\,S\longrightarrow\underline{\mathrm{Fibsc}}\,S,\quad
 \underline{\mathrm{Cofibsc}}_0\,S\longrightarrow\underline{\mathrm{Cofibsc}}\,S,
\]
% L13-14
étant pleinement fidèles et ayant comme images essentielles les sous-catégories en question.

% L15-20
Cette rétraction m'apparaît comme une chose tout à fait mystérieuse. \deleted{Pourtant, il semblerait que pour \unsure{le} \deleted{\unsure{dérivé}} dérivateur $\mathbb D$ (pas seulement HOT) qui provient \unsure{d'une|d'un} $(M,W)$ \deleted{avec $M$} provenant d'un \uline{triple de Quillen},} C'est lié au fait que pour le dérivateur très spécial $\mathbb D=\mathrm{HOT}$,
% L21-22
quand on a un $\xi\in\mathbb D(S)$, on en déduit \deleted{\unsure{une}} (à $\unsure{\mathrm{HOT}_{W/S}|\mathrm{HOT}_{/S}}$-équivalence près) un \deleted{$\unsure{\xi}$} \deleted{\unsure{type de}} \inserted{diagramme $X$ au dessus de $S$,} \deleted{\unsure{en quelque}}
% archive sheet number: 89
